% Proposed manuscript fragment. Retains the host edition's sans-serif text/math setup.
% Registry revision 20; SHA-256 cc9676276550560eddc941ee08874019d7566486be709a439c732b0996a1c1e8.
% Source-line provenance and outstanding documentary limits are in SOURCE-LEDGER.json.
\section{Frozen LINK Development Methods}

A predictor can lower its average error while making one finger group worse on
one recording date. The LINK comparison therefore keeps each date and output
visible, then asks which comparisons survive all their stated conditions.
The prospective protocol states the intended comparison. The retained
implementation specifies the operations; the completed-study report describes
the execution and its outcomes. Agreement among these records remains to be verified.

\subsection{Observations and Common Support}

The development design uses one macaque's four training dates, January 27 and 30
and February 4 and 5, 2020, and two validation dates, February 6 and 11, 2020.
The two measured targets are index-finger and combined middle/ring/small-finger
normalized positions. They have source-normalized position units, without a
conversion to radians, force or 23 separate joint angles. The input comprises
96 nonnegative threshold-crossing counts transformed by $\log(1+x)$ and 96
archived spiking-band-power features. Physical electrode contacts identify
columns; matching a contact across recordings does not identify a persistent
single neuron. These data do not supply simultaneous peripheral-nerve, spinal
or EMG recordings.

For a scored endpoint $t$, the input contains source frames $t-10,\ldots,t$ and
the auxiliary target contains the neural features at $t+1$. The position target
is at $t$. All twelve frames must belong to one source trial, and each adjacent
source-clock interval must lie in $[0.0185,0.0215]$ seconds. The nominal bin
interval is 20 milliseconds; original timestamp jitter is retained. No window
crosses a trial boundary, no interpolation fills a gap, and recurrent state
does not continue between windows. Every candidate, including the constant
and both ridge baselines, uses this same endpoint support and requires the
following neural frame. A position-only model consequently receives no extra
trial-end examples unavailable to the auxiliary model.

The first training date defines the canonical order of 96 physical contacts.
Both feature blocks are reordered together. Feature means and population
standard deviations are fitted in float64 on each unique trial-owned frame
once across the four training dates, rather than counting a frame again for
every overlapping window. Exactly zero standard deviations become one; there
is no other floor, clipping or validation-derived normalization. The transform
is calculated in float64 and cast to float32. The constant baseline is the
mean of unique trial-owned training positions. Position targets remain source
float32 values. Nonfinite measurements, unknown contacts or incompatible
schemas cause a failure rather than silent row removal or imputation.

The archived threshold-crossing unit attribute says volts despite count-valued
semantics; the comparison makes no voltage conversion. Archived power values
are retained without the author loader's unexplained factor of four. The
documented 32-millisecond interval also differs from the retained nominal
20-millisecond clock. Full author trial counts must agree with the trial stop
timestamps; uniform one-short time-series references are identified separately.
These distinctions matter because changing units or trial support would change
the experiment.

\subsection{Candidates, Initialization and Training Order}

The two recurrent candidates have 192 inputs, one gated recurrent-unit layer
with 128 hidden values, and two position outputs. The position head combines
the final hidden values with the mean neural feature vector over the eleven
input bins. A separate linear head predicts the 192 next-neural features.
Both candidates contain 169,090 parameters. The position-only objective uses
144,322; the 24,768-parameter neural head receives zero objective weight in that
candidate. The historical software name of the position head does not change
its target into velocity.

For position target $y$, position prediction $\widehat y$, next-neural target
$z$ and prediction $\widehat z$, training minimizes
\[
 L=\operatorname{mean}\bigl((\widehat y-y)^2\bigr)
   +\lambda\operatorname{mean}\bigl((\widehat z-z)^2\bigr),
 \qquad \lambda\in\{0,0.1\}.
\]
Each residual is squared before the mean is taken over the corresponding
minibatch and output elements. Otherwise, residuals of opposite signs could
cancel before the square and conceal prediction error. Training
uses Adam with learning rate $0.001$, averaging constants $(0.9,0.999)$,
stabilizer $10^{-8}$, zero weight decay, and disabled fused and foreach routes.
There is no dropout, autocast, TF32, gradient clipping or target clipping.
Minibatches contain 256 windows, including the final partial minibatch.
Nonfinite losses, gradients, updated parameters, optimizer states or predictions
fail the run. The reported study route uses four CPU mathematical threads and
one framework interoperation thread; that historical configuration is separate
from manuscript preparation and does not establish a speed improvement.

The fixed seeds are 20261007, 20261008 and 20261009. Each run resets Torch's
seed, constructs the same architecture, checks its initialized float32 tensors
against the seed's frozen initialization, and loads that initialization. The
intended pairing therefore includes both recurrent candidates and all three
input conditions for a given seed. A separate NumPy PCG64 generator starts
once per run from that seed. Training dates are concatenated in the declared
order above, with ascending scored source endpoints within each date. One
permutation of all concatenated training rows is drawn per epoch; every row
appears once. The implementation obtains matching schedules by resetting this
generator identically in each paired run. It saves each epoch's order hash and
generator state, rather than reading a shared bank of saved order arrays.

The run order is measured inputs, association control, then order control;
within a condition it is position-only, then position plus next-neural; within
a candidate it is the three seeds in their declared order. Each run retains
all 40 epochs. Restoring a completed epoch restores model, Adam, NumPy and Torch
state; an unacknowledged partial epoch is not permission to repeat updates.
The source specifies these safeguards. This documentary account does not
independently verify their use throughout the completed execution.

\subsection{What Each Control Changes}

The association control breaks the training alignment between neural evidence
and position. For each seed, the implementation creates one PCG64 generator
from \texttt{seed+1000000}, then advances that same stream across the four
training dates. It draws one permutation separately within each date. The
shuffle unit is an eligible whole-window row: its complete $11\times192$
neural window and associated 192-element next-neural target move together.
Position targets and the destination rows' date, trial, source-endpoint and
scoring identities stay fixed. This is neither a channel shuffle nor a
within-window bin shuffle, and it does not move whole trials as blocks.
The permutation is fixed for the run, shared by the two candidate types for
that seed, and is not redrawn each epoch. Its hash and accidental fixed-point
fraction are recorded.

Association-control models are fitted on these transformed training rows and
selected on unchanged measured validation windows and positions. Their final
comparison also uses unchanged validation support. Training losses refer to
the transformed training pairing; the main candidate and control scores refer
to the two validation dates. The control retains within-window history,
within-date distributions and possible common task structure. Fixed points
can retain some original pairings. It does not create a causal null.

The temporal-order control applies the same fixed permutation to the bins of
every training and validation window. In zero-based bin indices it reads
\[
 (9,8,7,6,5,4,3,2,1,0,10)
 \quad\hbox{in place of}\quad (0,1,2,3,4,5,6,7,8,9,10).
\]
Thus the ten earlier bins reverse and the current bin stays last. Position,
next-neural targets, feature identities and scoring rows remain unchanged.
The transformation uses no random draw. Each order-control run retains the
same seed initialization, epoch-order rule, optimizer and selection policy as
its measured-input counterpart. Reversal preserves the window's multiset of
values and is invertible; a recurrent model can adapt to it. The comparison
tests benefit from this declared order, not an assurance that all temporal
information has been removed. Both control conditions are excluded from the
best-model competition.

\subsection{Selection and Uncertainty}

Let $E_{dg}(A)$ be model $A$'s frame-mean squared position error for validation
date $d$ and finger group $g$. The primary score is
\[
 \overline E(A)=\frac14\sum_{d=1}^{2}\sum_{g=1}^{2}E_{dg}(A).
\]
Frames are weighted equally within a cell; the four date/group cells then have
equal weight. Each seed retains the epoch with the lowest primary score;
an exact tie takes the earliest epoch. Predictions from all three selected
seeds are averaged arithmetically with equal weights. Seed performance is
also retained separately, including its mean, sample standard deviation with
one degree-of-freedom correction, minimum and maximum. No favourable seed
replaces the ensemble.

The current-bin ridge uses only the final 192 features; history ridge flattens
all eleven bins oldest to newest in row-major order. Both fit float64 ridge
regression with unpenalized intercept and objective
$\sum_i\lVert y_i-X_iW-b\rVert^2+\alpha\lVert W\rVert_F^2$.
The frozen penalties are $0.01,0.1,1,10,100$. The lowest validation primary
score selects the penalty, with the smallest penalty winning an exact tie.
An exact tie between ridge types selects current-bin ridge. An exact tie
between measured recurrent ensembles selects position-only. Control models
do not enter either candidate-selection contest.

For every comparison, candidate-minus-reference squared errors are paired on
the same validation endpoints. The bootstrap initializes PCG64 with 20261007,
uses validation dates in declared order and sorted unique source trial IDs,
and draws 1,000 samples. Within each date, each sample draws the original
number of eligible trials with replacement. Every supported frame of a drawn
trial contributes with its multiplicity; variable trial lengths retain their
frame weighting. The same trial draws serve both models and both groups.
The generator advances from the first validation date to the second, and each
comparison restarts the same seed. The primary difference averages the four
cell differences. The 95-percent interval uses the 0.025 and 0.975 quantiles
with linear interpolation. Fewer than two supported trials per date prevents
this uncertainty calculation. The intervals are descriptive and conditional
on these dates and validation-selected models.

Ordinary $R^2$ uses the evaluation cell's observed mean:
\[
 R^2_{dg}(A)=1-
 \frac{\sum_i(\widehat y_{i,g}-y_{i,g})^2}
 {\sum_i(y_{i,g}-\overline y_{d,g})^2}.
\]
The separately labelled training-mean skill score substitutes the frozen
training mean in the denominator. A zero denominator produces an undefined
score, without an added epsilon, and cannot pass a rule requiring defined
$R^2$.

\subsection{Model Admission and the Five Benefit Comparisons}

Define the comparison conjunction
\[
 C(A,B)=
 [\overline E(A)<\overline E(B)]\ \land\
 [E_{dg}(A)\leq E_{dg}(B)\ \hbox{for every }d,g]\ \land\
 [U_{.95}(A-B)<0],
\]
where $U_{.95}$ is the upper endpoint of the paired two-sided 95-percent
interval just defined. Equality passes the per-cell no-harm condition, but
does not pass either strict improvement condition. Let $P$ and $N$ denote
the measured position-only and auxiliary ensembles, $R_c$ current-bin ridge,
$R_*$ selected best ridge, and $P_o,N_o$ their matched order controls.
The retained independent postfit checker evaluates these five comparisons separately:
\begin{enumerate}
 \item Position-only history benefit: $C(P,R_c)$.
 \item Position-only order-use benefit: $C(P,P_o)$.
 \item Auxiliary-model history benefit: $C(N,R_c)$.
 \item Auxiliary-model order-use benefit: $C(N,N_o)$.
 \item Auxiliary-objective benefit: $C(N,P)$.
\end{enumerate}
Each displayed conjunction includes its own lower primary error, no worse
date/group and negative paired upper bound. The history comparison asks
whether the recurrent candidate improves on the specified current-bin linear
reference. It does not isolate recurrent ordering from the mean-history input
or every other architectural difference. An order-use claim additionally
needs its matched reversal comparison; a history comparison alone does not
establish that claim.

The distinct implemented model-admission gate requires $C(A,R_*)$ and, in every
cell, defined ordinary $R^2(A)>0$ and $R^2(A)>R^2(A_s)$, where $A_s$ is the
matched association-control ensemble. A missing or nonfinite cell cannot pass.
This model gate must not be substituted for the five comparison rules.

The retained independent postfit checker implements the three-part conjunction
separately for every comparison above and records each Boolean condition. All
five decisions fail. Each has a lower average error, but each fails its
per-date/group no-harm condition; three have a paired upper bound below zero.
The separately recorded model-admission decisions also fail. This account
reuses those existing decisions without recomputing an interval or prediction.

That support is limited to the correctly negative report. The retained
order-control ledgers have empty control-provenance fields: source identity and
earlier accepted control evidence do not supply a new dynamic reversal check.
Nor does postfit agreement independently certify the shared optimizer, every
new checkpoint forward calculation, restored random state or complete update
trajectory. Those requirements remain open. A favourable average cannot
replace a failed condition.

\subsection{Completion and Causal Limits}

The inventory records completion of 18 runs with 40 saved epochs each, for 720 saved epochs. That completion and the later postfit checks support the unchanged negative report; they do not verify every optimizer calculation, checkpoint or update, or establish whole-study scientific acceptance. The prospective protocol's \texttt{execution\_started=false} records its state when frozen, before the later study ran. Completion admits neither a model nor reserved confirmation.

The saved inventory records that the service was terminal and inactive, with exit status zero and no automatic restart. That record establishes custody of the files; it does not, by itself, bind them to the complete execution command, loaded modules and later decision records. Those missing links limit what can be reconstructed about execution. They do not explain away the failed development comparisons.

February 22 and 24, 2020, remain closed later-day dates; June 21 and 22, 2023,
remain closed drift dates. The later task difference also limits a pure drift
interpretation. No reserved result is inferred from development validation.

The history is causal only with respect to the supplied source indices:
position prediction uses indices no later than $t$, while next-neural learning
uses $t+1$ as a training target. Full archive extraction and filtering support
has not established causal online acquisition or physiological delay. The
study concerns predictive associations between recorded features and two
position coordinates. It does not demonstrate nerve transmission, spinal
externalization, sensory delivery, biological teaching, retained skill or
full-body neural translation.
